%! AP Statistics — Unit 1, Lesson 1.9 — Warm-Up ANSWER KEY
\documentclass[10pt]{article}
\usepackage{apstats-article}
\usepackage{apstats-key}

%=============================================================================
\begin{document}

\pageheader{Unit 1, Lesson 1.9}{Warm-Up --- Answer Key}
\begin{tabularx}{\linewidth}{X X X}
Name:\;\ans{Key} & Date:\; & Period:\;
\end{tabularx}

\vspace{0.16in}

\noindent In Lessons 1.7--1.8 you built and read boxplots for a single group.
Today you will compare two groups on the \textit{same} display.
The parallel boxplots below show \textbf{ACT composite scores} for two groups:
those who did \textit{not} take a test-prep course (\textbf{No Prep})
and those who did (\textbf{Test Prep}).

\vspace{0.18in}

\begin{center}
\begin{tikzpicture}[scale=1.0]
  \draw[thick] (0.0, 0) -- (6.80, 0);
  \foreach \v/\x in {14/0.00, 16/0.60, 18/1.20, 20/1.80, 22/2.40,
                      24/3.00, 26/3.60, 28/4.20, 30/4.80, 32/5.40, 34/6.00, 36/6.60} {
    \draw (\x, -0.12) -- (\x, 0.12);
    \node[below, font=\small] at (\x, -0.12) {\v};
  }
  \node[below, font=\small\itshape] at (3.30, -0.56) {ACT composite score};
  \node[left, font=\small] at (0.0, 1.10) {No Prep};
  \draw[thick, navy] (0.30, 1.10) -- (1.50, 1.10);
  \draw[thick, navy] (1.50, 0.85) rectangle (3.30, 1.35);
  \draw[thick, navy] (2.40, 0.85) -- (2.40, 1.35);
  \draw[thick, navy] (3.30, 1.10) -- (4.50, 1.10);
  \draw[thick, navy] (0.30, 0.95) -- (0.30, 1.25);
  \draw[thick, navy] (4.50, 0.95) -- (4.50, 1.25);
  \node[left, font=\small] at (0.0, 0.35) {Test Prep};
  \draw[thick, navy] (1.80, 0.35) -- (3.00, 0.35);
  \draw[thick, navy] (3.00, 0.10) rectangle (4.80, 0.60);
  \draw[thick, navy] (3.90, 0.10) -- (3.90, 0.60);
  \draw[thick, navy] (4.80, 0.35) -- (6.30, 0.35);
  \draw[thick, navy] (1.80, 0.20) -- (1.80, 0.50);
  \draw[thick, navy] (6.30, 0.20) -- (6.30, 0.50);
\end{tikzpicture}
\end{center}

\vspace{0.20in}

\begin{enumerate}[leftmargin=1.5em, itemsep=8pt]

  \item \ans{No Prep median $= 22$; Test Prep median $= 27$.
        The Test Prep group scored higher on average.}

  \item \ans{No Prep: $Q_1 = 19$, $Q_3 = 25$, IQR $= 6$ points.
        Test Prep: $Q_1 = 24$, $Q_3 = 30$, IQR $= 6$ points.
        The two groups have equal IQRs, so their middle 50\% scores
        are equally spread.}

  \item \ans{Yes, the distributions overlap from 20 to 29 --- scores in that
        range appear in both groups.}

  \item \ans{The median ACT score for the Test Prep group (27) was 5 points
        higher than the median for the No Prep group (22), suggesting that
        students who completed the test-prep course tended to score higher.}

\end{enumerate}

\vspace{0.14in}

\begin{tcolorbox}[colback=greenbg, colframe=greenacc, boxrule=0.8pt, arc=2mm,
    left=4mm, right=4mm, top=3mm, bottom=3mm]
\small
\begin{itemize}[leftmargin=1.3em, itemsep=5pt]
  \item Today you will learn how to use the SOCS framework to compare
        \textit{two groups} on the same display --- not just describe one.
  \item You will also practice a back-to-back stemplot, a second display
        for comparing small datasets.
  \item Key focus: AP scoring requires \textbf{explicit} comparative
        language. Describing each group separately does not earn credit.
\end{itemize}
\end{tcolorbox}

\begin{teachernote}
\textbf{Grading notes (4 pts total):}
\begin{itemize}[itemsep=2pt]
  \item Q1: 1 pt --- both medians correct (22, 27) AND names Test Prep as higher.
  \item Q2: 1 pt --- all four quartile values correct and IQRs = 6; correct spread conclusion.
  \item Q3: 1 pt --- identifies overlap AND gives approximate shared range (any reasonable range from 20--29).
  \item Q4: 1 pt --- uses comparative language with both group names AND includes context (ACT scores).
\end{itemize}
\end{teachernote}

\end{document}
